\subsection{本阶段做得较好的地方}

这一阶段最有价值的变化，是分工开始落到模型接口，而不只是认领几个算法名称。队长
输出位置和速度，组员1把这些量写进速度约束，组员2维护参数表和误差口径，三人的结果
可以在同一张变量表里接上。几个简单基线也保留了下来：RK4对照Euler，MILP对照小规模
枚举，粒子群对照遍历，树模型对照线性或Logistic。读A053时，笔记也没有停在目录层面，其中
“定长递推--物理选根--第一次碰撞--回退求精--分区状态”这条顺序已经整理成一份可逐项核对
代码的笔记。

